\subsection{Sparse Temporal Sampling: 16 vs.\ 8 vs.\ 4 Frames (Frozen DINOv2, SSv2)}
\label{exp:sparse-temporal-sampling}

\paragraph{What we tried} Scored plain frozen DINOv2 ViT-S/14 CLS features on SSv2
validation at three evenly-sampled frame counts ($T=16$, 8, 4) --- no training, no
regularizer, just subsampling the same clips more sparsely --- to check whether the
metrics used throughout this log are comparable across different $T$, since most prior
entries fix $T=16$. Scored with the same full metric suite from
Section~\ref{sec:methodology} (Main and Additional).

\paragraph{Result}

\begin{table}[H]
\centering
\caption{Main metrics, frozen DINOv2 CLS, SSv2, at three sampling densities.}
\label{tab:sparse-sampling-main}
\begin{tabular}{lccc}
\toprule
Metric & 16-frame & 8-frame & 4-frame \\
\midrule
\multicolumn{4}{l}{\textit{Whole-dataset collapse and spread}} \\
effective\_rank\_fraction & 0.8120 & 0.8116 & 0.8105 \\
$W$ (within-clip mean pairwise L2) & 15.3090 & 15.8230 & 16.6857 \\
$B$ (across-clip mean pairwise L2) & 31.1109 & 31.0909 & 31.0612 \\
$B-W$ (gap) & 15.8019 & 15.2679 & 14.3755 \\
\midrule
\multicolumn{4}{l}{\textit{Trajectory geometry}} \\
trajectory\_straightness ($z$) & 0.1348 & \good{0.2277}\textsuperscript{d} & \good{0.4241}\textsuperscript{d} \\
mean\_velocity\_cosine ($z$) & -0.2636 & \bad{-0.2764} & \bad{-0.3106} \\
start\_end\_displacement ($z$) & 0.8001 & 0.7948 & 0.7981 \\
direction\_reversal\_fraction ($x$) & 0.8941 & \bad{0.9093} & \bad{0.9322} \\
\midrule
\multicolumn{4}{l}{\textit{Temporal order}} \\
time\_progress\_spearman ($z$) & 0.9040 & \good{0.9247}\textsuperscript{d} & \good{0.9699}\textsuperscript{d} \\
distance\_to\_goal\_spearman ($x$) & 0.6110 & \bad{0.5808} & \bad{0.5143} \\
adjacent\_temporal\_nn\_fraction ($z$) & 0.7873 & 0.8014 & 0.7984 \\
\midrule
\multicolumn{4}{l}{\textit{Cross-clip temporal consistency}} \\
tcc\_kendall\_tau (mean) & 0.0661 & 0.0762 & 0.1258 \\
\midrule
\multicolumn{4}{l}{\textit{Semantic information}} \\
chirality.mean\_accuracy ($x$) & 0.7993 & 0.7784 & \bad{0.7455} \\
content.knn.balanced\_acc & 0.2050 & 0.2063 & 0.2044 \\
content.linear.balanced\_acc & 0.2521 & 0.2436 & 0.2560 \\
change.knn.balanced\_acc & 0.3110 & 0.3048 & 0.0455\textsuperscript{a} \\
change.linear.balanced\_acc & 0.3681 & 0.3588 & 0.0455\textsuperscript{a} \\
content+change.knn.balanced\_acc & 0.3434 & 0.3400 & 0.2044 \\
content+change.linear.balanced\_acc & 0.5242 & 0.4912 & 0.2569 \\
\bottomrule
\end{tabular}

\vspace{0.3em}
{\footnotesize \textsuperscript{a}\texttt{change} is defined as mean(last \texttt{edge\_frames})
$-$ mean(first \texttt{edge\_frames}), default \texttt{edge\_frames=4}. At $T=4$ the two
windows are identical (all 4 frames), so \texttt{change}$\equiv 0$ for every clip and
both probes collapse exactly to \texttt{chance.uniform} (1/22). Not evidence about the
encoder --- a mechanical consequence of the default edge width meeting $T=4$ --- and a
reason \texttt{--edge-frames} would need to shrink for any future $T<8$ comparison
involving the \texttt{change} vector specifically. \quad
\textsuperscript{d}Green here is \underline{not} a real improvement --- see Inference
below: fewer waypoints mechanically inflates both the path-length/net-displacement ratio
(\texttt{trajectory\_straightness} and relatives) and Spearman's rank correlation
(\texttt{time\_progress\_spearman}, fewer points $=$ fewer possible rank orderings to get
wrong), independent of trajectory quality. Do not read these cells the way \good{green}
is read everywhere else in this document.}
\end{table}

\begin{table}[H]
\centering
\caption{Additional metrics, frozen DINOv2 CLS, SSv2 --- same three sampling densities as
Table~\ref{tab:sparse-sampling-main}.}
\label{tab:sparse-sampling-additional}
\begin{tabular}{lccc}
\toprule
Metric & 16-frame & 8-frame & 4-frame \\
\midrule
\multicolumn{4}{l}{\textit{Trajectory geometry and raw variants}} \\
mean\_chord\_alignment ($z$) & 0.1146 & \good{0.1998}\textsuperscript{d} & \good{0.3881}\textsuperscript{d} \\
trajectory\_straightness\_raw ($x$) & 0.1358 & \good{0.2290}\textsuperscript{d} & \good{0.4260}\textsuperscript{d} \\
mean\_velocity\_cosine\_raw ($x$) & -0.2627 & \bad{-0.2750} & \bad{-0.3077} \\
mean\_chord\_alignment\_raw ($x$) & 0.1154 & \good{0.2010}\textsuperscript{d} & \good{0.3900}\textsuperscript{d} \\
\midrule
\multicolumn{4}{l}{\textit{Raw embedding magnitude}} \\
norm\_growth ($x$) & 0.9933 & 0.9952 & 0.9977 \\
norm\_std\_within\_clip ($x$) & 0.5421 & 0.5319 & 0.4978 \\
mean\_norm ($x$) & 23.653 & 23.645 & 23.631 \\
\midrule
\multicolumn{4}{l}{\textit{Change-direction consistency}} \\
change\_direction.gap & 0.0386 & 0.0338 & 0.0000\textsuperscript{a} \\
\bottomrule
\end{tabular}

\vspace{0.3em}
{\footnotesize \textsuperscript{a}see Table~\ref{tab:sparse-sampling-main}'s footnote.
\quad \textsuperscript{d}same caveat as Table~\ref{tab:sparse-sampling-main}: mechanical
inflation from fewer waypoints, not a real straightness improvement.}
\end{table}

\paragraph{Inference}

\textbf{This closes the log with a methodology check rather than another checkpoint:
every comparison in every experiment above fixes $T=16$, so before trusting any of those
comparisons, it's worth confirming the metrics themselves don't silently depend on frame
count.} They do, for some of them --- which is exactly why this entry comes last: it's a
caveat on how to read everything before it, not a finding that stands on its own the way
the checkpoints above do.

\textbf{\texttt{trajectory\_straightness} (and its raw/chord-alignment relatives) AND
\texttt{time\_progress\_spearman} are \emph{not} invariant to frame count --- both are
green for 8-/4-frame in the table above, and both are explicitly flagged rather than
trusted.} Comparing either across different $T$ is comparing metric resolution, not real
trajectory quality: straightness rises from 0.1348 ($T{=}16$) to 0.2277 ($T{=}8$) to
0.4241 ($T{=}4$), and \texttt{time\_progress\_spearman} rises the same way (0.9040 $\to$
0.9247 $\to$ 0.9699), on the \emph{exact same encoder and clips} --- nothing trained or
changed except sampling density. Fewer waypoints mechanically gives less opportunity for
the path-length/net-displacement ratio to register zig-zag, and fewer points give
Spearman fewer possible rank orderings to get wrong. \texttt{trajectory\_straightness} is
the same quantity that turned green for every collapsed checkpoint earlier in this log
(Section~\ref{exp:window-vicreg-ablation}, Section~\ref{exp:regularizer-comparison-k400}'s
\texttt{noreg}) --- there the cause was collapse; here, on a completely healthy encoder,
the identical metric inflates for a totally different, purely mechanical reason. Two
different failure modes produce the same-looking green cell, which is the strongest
reason not to trust \texttt{trajectory\_straightness} --- or, per this experiment,
\texttt{time\_progress\_spearman} --- alone anywhere in this document.

\textbf{\texttt{distance\_to\_goal\_spearman} moves the \emph{opposite} direction under
sparser sampling, and is red rather than green --- it argues for being the metric to
trust here, not the one to discount.} It falls as $T$ drops (0.6110 $\to$ 0.5808, red
$\to$ 0.5143, red), and \texttt{chirality.mean\_accuracy} falls too (0.7993 $\to$ 0.7784,
uncolored $\to$ 0.7455, red) --- both getting worse with sparser sampling, the opposite
trend from the two green-but-flagged metrics above. This makes sense for genuine reasons:
distance-to-goal needs enough intermediate points to show a smooth monotonic approach,
and telling forward from reverse needs enough temporal resolution to detect the
direction of change. Both are plausibly measuring real information loss as frames get
sparser, unlike the mechanical inflation above --- so of the two "temporal order"
metrics used throughout this log, \texttt{distance\_to\_goal\_spearman} is the one that
behaves like a genuine quality signal under $T$, and \texttt{time\_progress\_spearman} is
the one that doesn't.

\textbf{\texttt{start\_end\_displacement} and the content probes are correctly invariant,
which is a useful sanity check on the pipeline itself.} \texttt{start\_end\_displacement}
only depends on the first and last frame, so it stays flat (0.8001/0.7948/0.7981)
regardless of intermediate sampling, as it should. \texttt{content.knn.balanced\_acc} and
\texttt{content.linear.balanced\_acc} are similarly flat (0.2044--0.2063 /
0.2436--0.2560) --- semantic content is a property of the frames themselves, not of how
many are used to build the trajectory, and this confirms the metric pipeline reflects
that correctly rather than leaking a $T$-dependence it shouldn't have.

\textbf{Practical implication for reading the rest of this log.} Every
\texttt{trajectory\_straightness}-family and \texttt{time\_progress\_spearman} number
elsewhere in this document is implicitly compared only against other $T{=}16$ runs,
which every other entry in this log is. If a future experiment needs a different frame
count, those numbers are not directly comparable to this log's existing tables without
accounting for $T$ --- \texttt{distance\_to\_goal\_spearman}, \texttt{chirality.mean\_accuracy},
and the content probes would be the safer metrics to lean on for any such cross-$T$
comparison.
